\subsection{分次同态的分次模}

我们在本号中假设幺半群 \(\Delta\) 是一个群。设 A 是类型为 \(\Delta\) 的分次环，M、N 是两个类型为 \(\Delta\) 的分次左（例如）A-模。设 H \(_{\lambda}\) 表示由 M 到 N 的次数为 \(\lambda\) 的分次同态构成的加法群（第 2 号）；在 M 到 N 的所有同态（带非分次 A-模结构）的加法群 Hom \(_{A}\) (M, N) 中，对 \(\lambda \in \Delta\) 的诸 H \(_{\lambda}\) 之和是直和。因为，若有一个关系 \(\sum_{\lambda} u_{\lambda} = 0\) ，其中 \(u_{\lambda} \in \mathrm{H}_{\lambda}\) （对所有 \(\lambda\) ），则由此得出 \(\sum_{\lambda} u_{\lambda}(x_{\mu}) = 0\) （对所有 \(\mu\) 及所有 \(x_{\mu} \in \mathrm{M}_{\mu}\) ）。由于 \(\Delta\) 的元素都可消去， \(u_{\lambda}(x_{\mu})\) 是 \(\sum_{\lambda} u_{\lambda}(x_{\mu})\) 的次数为 \(\lambda + \mu\) 的齐次分量；因此 \(u_{\lambda}(x_{\mu}) = 0\) （对每个有序对 ( \(\mu, \lambda\) ) 及每个 \(x_{\mu} \in \mathrm{M}_{\mu}\) ），这蕴含 \(u_{\lambda} = 0\) （对所有 \(\lambda \in \Delta\) ）。我们（在本段中）用 Homgr \(_{A}\) (M, N) 表示 Hom \(_{A}\) (M, N) 的加法子群，即所有 H \(_{\lambda}\) 的和，并称之为从 M 到 N 的分次 A-模同态的加法群。设 C 为 A 的中心，它是一个分次子环（第 3 号，命题 5 的推论）；对于 \(\mathrm{Hom}_{\mathbf{A}}(\mathbf{M},\mathbf{N})\) 上的典范 C-模结构（§ 1，第 14 号，注记 1）， \(\mathrm{Homgr}_{\mathbf{A}}(\mathbf{M},\mathbf{N})\) 是一个子模，且分次 \((\mathbf{H}_{\lambda})\) 与 C-模结构相容：因为，若 \(c_{\nu}\in \mathbf{C}\cap \mathbf{A}_{\nu}\) , \(x_{\mu}\in \mathbf{N}_{\mu}\) 和 \(u_{\lambda}\in \mathbf{H}_{\lambda}\) ，则由定义 \((c_{\nu}u_{\lambda})(x_{\mu}) = c_{\nu}.u_{\lambda}(x_{\mu})\subset \mathbf{N}_{\lambda +\mu +\nu}\) ，因此 \(c_{\nu}u_{\lambda}\in \mathbf{H}_{\lambda +\nu}\) .

设 \(\mathbf{M}'\) 和 \(\mathbf{N}'\) 是两个类型为 \(\Delta\) 的分次左 A-模，且 \(u':\mathbf{M}'\to\mathbf{M}\) , \(v':\mathbf{N}\to\mathbf{N}'\) 是次数分别为 \(\alpha\) 和 \(\beta\) 的分次同态。那么立即看出 \(\operatorname{Hom}(u',v')\colon w\mapsto v'\circ w\circ u'\) 把 \(\operatorname{Homgr}_{\mathbf{A}}(\mathbf{M},\mathbf{N})\) 映到 \(\operatorname{Homgr}_{\mathbf{A}}(\mathbf{M}',\mathbf{N}')\) 中，且它到 \(\operatorname{Homgr}_{\mathbf{A}}(\mathbf{M},\mathbf{N})\) 的限制是到 \(\operatorname{Homgr}_{\mathbf{A}}(\mathbf{M}',\mathbf{N}')\) 中的次数为 \(\alpha+\beta\) 的分次同态。

特别地， \(\operatorname{Homgr}_{\mathbf{A}}(\mathbf{M}, \mathbf{M})\) 是 \(\operatorname{End}_{\mathbf{A}}(\mathbf{M})\) 的一个分次子环，记作 \(\operatorname{Endgr}_{\mathbf{A}}(\mathbf{M})\) 。

注. 若 M 和 N 是分次左 A-模， \(\operatorname{Homgr}_{\mathbf{A}}(\mathbf{M}, \mathbf{N})\) 一般不同于 \(\operatorname{Hom}_{\mathbf{A}}(\mathbf{M}, \mathbf{N})\) 。然而当 M 是有限生成的 A-模时，这两个集合相等。因为此时 M 由有限个齐次元素 \(x_{i}\) （ \(1 \leqslant i \leqslant n\) ）生成；设 \(d(i)\) 是 \(x_{i}\) 的次数；令 \(u \in \operatorname{Hom}_{\mathbf{A}}(\mathbf{M}, \mathbf{N})\) ，且对所有 \(\lambda \in \Delta\) ，设 \(z_{i,\lambda}\) 表示 \(u(x_{i})\) 的次数为 \(\lambda + d(i)\) 的齐次分量。我们证明存在一个同态 \(u_{\lambda}: \mathbf{M} \to \mathbf{N}\) ，使得 \(u_{\lambda}(x_{i}) = z_{i,\lambda}\) （对所有 \(i\) ）。只需证明若 \(\sum_{i} a_{i} x_{i} = 0\) ，其中 \(a_{i} \in \mathbf{A}\) （对于 \(1 \leqslant i \leqslant n\) ），则 \(\sum_{i} a_{i} z_{i,\lambda} = 0\) （对所有 \(\lambda \in \Delta\) ）（§1, 第 7 号, 注）。可以假设每个 \(a_{i}\) 都是次数为 \(d'(i)\) 的齐次元素，使得 \(d(i) + d'(i) = \mu\) （对所有 \(i\) ）（第 3 号, 注 1）；则 \(\sum_{i} a_{i} u(x_{i}) = 0\) ；在左边取次数为 \(\lambda + \mu\) 的齐次分量，我们得到 \(\sum_{i} a_{i} z_{i,\lambda} = 0\) ，由此得出同态 \(u_{\lambda}\) 的存在性；此外显然 \(u_{\lambda}\) 是次数为 \(\lambda\) 的分次同态。最后， \(u_{\lambda} = 0\) （除 \(\lambda\) 的有限多个值外），且按定义 \(u = \sum_{\lambda} u_{\lambda}\) ，这就证明了我们的断言。

特别地， \(\operatorname{Homgr}_{\mathbf{A}}(\mathbf{A}_s, \mathbf{M}) = \operatorname{Hom}_{\mathbf{A}}(\mathbf{A}_s, \mathbf{M})\) （对每个分次左 A-模 M）；此外 \(\operatorname{Hom}_{\mathbf{A}}(\mathbf{A}_s, \mathbf{M})\) 具有一个分次左 A-模结构（而不仅仅是分次 C-模结构），且显然在此结构下，M 到 \(\operatorname{Hom}_{\mathbf{A}}(\mathbf{A}_s, \mathbf{M})\) 中的典范映射（ \(\S 1\) , 第 14 号, 注 2）是一个分次 A-模同构。

类似地， \(\operatorname{Homgr}_{\mathbf{A}}(\mathbf{M}, \mathbf{A}_s)\) 具有一个分次右 A-模结构（而不仅仅是分次 C-模结构）；它称为分次 A-模 M 的分次对偶，记作 \(\mathbf{M}^{*\mathbf{gr}}\) ，或当不会引起混淆时简记作 \(\mathbf{M}^*\) 。如果 \(u: \mathbf{M} \to \mathbf{N}\) 是次数为 \(\delta\) 的分次同态，则由上述可知，到 \(\mathbf{N}^{*\mathbf{gr}}\) 上的 \(^t u = \operatorname{Hom}(u, 1_{\mathbf{A}_s})\) 的限制，是分次对偶 \(\mathbf{N}^{*\mathbf{gr}}\) 到分次对偶 \(\mathbf{M}^{*\mathbf{gr}}\) 中的一个次数为 \(\delta\) 的分次同态，称为 \(u\) 的分次转置。

我们有时在分次对偶 \(M^{*gr}\) 上考虑借助同构 \(\lambda\mapsto-\lambda\) （ \(\Delta\) 的）从上述分次导出的分次（第 1 号，例 2），使得次数为 \(\lambda\) 的齐次元素（ \(M^{*gr}\) 中的）是 M 上次数为 \(-\lambda\) 的分次线性形式（当 A 具有平凡分次时，这些是在 \(M_{\mu}\) （指标 \(\mu\neq\lambda\) 的）上为零的线性形式）。那么，若 \(u:M\to N\) 是次数为 \(\delta\) 的分次同态，则 \(u'\) 成为次数为 \(-\delta\) 的分次同态。

设 A 是交换的且类型为 \(\Delta\) 的分次环，且设 M、N、P、Q 是四个类型为 \(\Delta\) 的分次 A-模。则存在次数为 0 的典范分次同态

\begin{enumerate}
\def\labelenumi{(\arabic{enumi})}
\setcounter{enumi}{2}
\tightlist
\item
\end{enumerate}

\[
\operatorname{Homgr} _ {\mathbf {A}} (\mathbf {M}, \operatorname{Homgr} _ {\mathbf {A}} (\mathbf {N}, \mathbf {P})) \rightarrow \operatorname{Homgr} _ {\mathbf {A}} (\mathbf {M} \otimes_ {\mathbf {A}} \mathbf {N}, \mathbf {P})\tag{4}
\]

\[
\operatorname{Homgr} _ {\mathbf {A}} (\mathbf {M}, \mathbf {N}) \otimes_ {\mathbf {A}} \mathrm{P} \rightarrow \operatorname{Homgr} _ {\mathbf {A}} (\mathbf {M}, \mathbf {N} \otimes_ {\mathbf {A}} \mathrm{P})
\]

\[
\operatorname{Homgr} _ {\mathbf {A}} (\mathbf {M}, \mathbf {P}) \otimes_ {\mathbf {A}} \operatorname{Homgr} _ {\mathbf {A}} (\mathbf {N}, \mathbf {Q}) \rightarrow \operatorname{Homgr} _ {\mathbf {A}} (\mathbf {M} \otimes_ {\mathbf {A}} \mathbf {N}, \mathbf {P} \otimes_ {\mathbf {A}} \mathbf {Q}) \tag {5}
\]

（张量积赋予第 5 号中所定义的分次），它们由限制 §4, 第 1、2、4 号中所定义的典范同态而得到；事实上，若 \(u: \mathbf{M} \to \operatorname{Homgr}_{\mathbf{A}}(\mathbf{N}, \mathbf{P})\) 是次数为 \(\delta\) 的分次同态，那么对所有 \(x \in \mathbf{M}_{\lambda}\) , \(u(x)\) 是 \(\mathbf{N} \to \mathbf{P}\) 的次数为 \(\delta + \lambda\) 的分次同态，因此对于 \(y \in \mathbf{N}_{\mu}\) , \(u(x)(y) \in \mathbf{P}_{\delta + \lambda + \mu}\) ；若 \(v: \mathbf{M} \otimes_{\mathbf{A}} \mathbf{N} \to \mathbf{P}\) 典范地对应于 \(u\) ，则于是看出 \(v\) 是次数为 \(\delta\) 的分次同态，由此得出关于 (3) 的断言；此外还可看出这个同态是双射。对 (4) 和 (5) 的论证类似。

特别地，若 \(\mathbf{P} = \mathbf{Q} = \mathbf{A}\) 在(5)中，则存在一个次数为0的典范分次同态

\[
\mathbf {M} ^ {* \mathrm{gr}} \otimes_ {\mathbf {A}} \mathbf {N} ^ {* \mathrm{gr}} \rightarrow (\mathbf {M} \otimes_ {\mathbf {A}} \mathbf {N}) ^ {* \mathrm{gr}}.\tag{6}
\]

